\begin{Proposition}{整环上的多项式}\label{prop:polynomials-over-domain}
  若 $R$ 是整环, 则 $R[x_{1},\ldots,x_{n}]$ 是整环.
\end{Proposition}

\begin{Proof}
  我们对变量的个数作归纳.

  \textbf{(归纳起点)} 设 $f,g\in R[x_{1}]$ 非零. 我们必须证明 $fg\ne0$. 写成
  \[
    f=a_{d}x_{1}^{d}+a_{d-1}x_{1}^{d-1}+\cdots+a_{0}
    \quad\text{且}\quad
    g=b_{e}x_{1}^{e}+b_{e-1}x_{1}^{e-1}+\cdots+b_{0},
  \]
  其中 $a_{d}\ne0$ 且 $b_{e}\ne0$. 根据 $R[x_{1}]$ 中乘法的定义,
  \[
    fg=(a_{d}b_{e})x_{1}^{d+e}
    +(a_{d}b_{e-1}+a_{d-1}b_{e})x_{1}^{d+e-1}+\cdots+a_{0}b_{0}.
  \]
  由于 $R$ 是整环, $a_{d}b_{e}\ne0$. 因为 $fg$ 至少有一个非零系数, 所以 $fg\ne0$.

  \textbf{(归纳步骤)} 假设 $S=R[x_{1},\ldots,x_{n-1}]$ 是整环; 我们必须证明 $R[x_{1},\ldots,x_{n}]$ 是整环. 由 Proposition~\ref{prop:polynomial-ring-recursion},
  \[
    R[x_{1},\ldots,x_{n}]=S[x_{n}].
  \]
  由于 $S$ 是整环, 归纳起点中的论证立即表明 $S[x_{n}]$, 从而 $R[x_{1},\ldots,x_{n}]$, 是整环.
\end{Proof}

Proposition~\ref{prop:polynomials-over-domain} 的证明中出现的数 $d$ 和 $e$ 是多项式 $f$ 和 $g$ 的重要属性, 称为它们的\textbf{次数}. 下一个定义把次数的概念推广到具有任意个变量的多项式.

\begin{Definition}{单项式和多项式的次数}\label{def:polynomial-degree}
  单项式 $x^{\alpha}=x_{1}^{\alpha_{1}}\cdots x_{n}^{\alpha_{n}}$ 的\textbf{次数}为
  \[
    \deg(x^{\alpha})=\alpha_{1}+\cdots+\alpha_{n}\in\mathbb{N}.
  \]
  非零多项式 $f=\sum a_{\alpha}x^{\alpha}\in R[x_{1},\ldots,x_{n}]$ 的\textbf{次数}为
  \[
    \deg(f)=\max\{\deg(x^{\alpha})\mid a_{\alpha}\ne0\}.
  \]
\end{Definition}

\begin{Example}{次数}\label{ex:monomial-and-polynomial-degrees}
  单项式 $x^{2}yz$, $z^{4}$, $x$ 和 $1$ 的次数分别为 $4$, $4$, $1$ 和 $0$, 并且
  \[
    \deg(x^{2}yz+z^{4}+x+1)=4.
  \]
\end{Example}

读者可能已经注意到, 我们没有给零多项式指定次数, 这是有意为之. 次数最有用的性质之一是下一个结果所述的\textbf{可加性}; 若选取任意 $\deg(0)\in\mathbb{N}$, 这个性质都不成立.
